\scanpage{95}
\deleted{\mspage{\unsure{24}}}
est la carte $P$-pondérée \inserted{orientée} cherchée.

Pour avoir un foncteur quasi inverse de (101), on procède de façon analogue, en utilisant une orientation de $P$. (Sinon, on utilise l'ens $\underline\omega$ pour tordre le corps $\mathbb C$, et le remplacer par $\mathbb C(\omega)$, et on considère \deleted{la sphère} $\mathbb C(\omega)$-variété $\Sigma_P(\mathbb C(\omega))$ au lieu de $\Sigma_P(\mathbb C)$, qui aura un point canonique $P_+$ pouvant servir de pt base). On \deleted{passe} interprète $\Pi_P^\tau$ comme $\Pi_1(\Sigma_P^*,\{1,\tau\};P_\omega)$, et on procède comme p.~9 (dans le cas $n=3$, dont on a ici une généralisation \ldots).

Sûrement on doit pouvoir généraliser aussi au cas quelconque, les considérations du § 5 (p.~$12'$--14), qui précisent (en relation avec \deleted{\unsure{les}} les iso. (97) (99)) les relations entre cartes $P$-pondérées, et cartes $n$-gonales (i.e. dont toutes les faces sont $n$-gonales \ldots). Je ne vais pas détailler ce point ici.

\margin{94}
